\chapter*{Kurzzusammenfassung}

Robuste Optimierung soll Optimalität erreichen, während
sämtliche Nebenbedingungen erfüllt sind
–  trotz Störungen, Parameterabweichungen oder nicht modellierter Effekte.
In \emph{Model Predictive Control} (MPC) sind
eine schnelle Berechnung und Robustheit entscheidend.
Durch die Konformalisierung einer Quantilsregression (CQR) können
Neuronale Netze (NNs) Quantile von Verteilungen
lernen und die rechenintensive 
Integration im MPC-Algorithmus ersetzen.
Für einen lastflexiblen Festbettreaktor ist gezeigt, wie NNs mit CQR 
die Unsicherheit in der Temperatur quantifizieren.
Dadurch ermöglichen sie die vollständige Einhaltung der Maximaltemperatur 
während der MPC-Steuerung.
Zudem wird eine \emph{physics-constrained} (PC) Aktivierungsfunktion präsentiert, 
die NN-Ausgaben während Inferenz und Gradientenbildung
in einen definierten Raum zwingt. 
Durch den korrigierten Gradienten wird das Training beschleunigt.
Außerdem reduzieren PCNNs den 
MPC-Stellaufwand hier um \SI{8.5}{\%} und führen zu \SI{21}{\%} kürzeren Rechenzeiten,
was einer Glättung der Zielfunktion zugeschrieben wird.



\chapter*{Abstract}

Robust optimization aims to maintain high performance while satisfying constraints under perturbations,
parameter offsets, or unmodeled effects. This is especially 
demanding in model predictive control (MPC), where fast computation and robustness are crucial.
Neural networks (NNs) combined with conformalized quantile regression (CQR) can learn confidence 
intervals and replace expensive integration and state estimation in MPC.
For a load-flexible fixed-bed reactor, CQR-based NNs quantify temperature uncertainty
during recursive simulation and ensure zero violations of the maximum temperature constraint in MPC.
A physics-constrained (PC) activation function is introduced to map NN outputs into a 
feasible space during both forward and backward passes. This accelerates training by 
correcting NN gradients. In this case, PC-based NNs reduce MPC control effort by \SI{8.5}{\%}
and enable faster 
solution times by \SI{21}{\%}, assuming a smoother objective landscape.
